\section{Thiết kế Edge Runtime}

Luồng xử lý của Edge Runtime được thể hiện trong
Hình~\ref{fig:edge-runtime-design}.

\begin{figure}[H]
\centering
\begin{tikzpicture}[node distance=0.55cm]
  \node[flowbox] (camera) {Camera};
  \node[flowbox, below=of camera] (capture) {Frame Capture};
  \node[flowbox, below=of capture] (pre) {Preprocess};
  \node[flowbox, below=of pre] (trt) {TensorRT};
  \node[flowbox, below=of trt] (post) {Postprocess};
  \node[flowbox, below=of post] (det) {Detection};
  \draw[flowarrow] (camera) -- (capture);
  \draw[flowarrow] (capture) -- (pre);
  \draw[flowarrow] (pre) -- (trt);
  \draw[flowarrow] (trt) -- (post);
  \draw[flowarrow] (post) -- (det);
\end{tikzpicture}
\caption{Thiết kế pipeline Edge Runtime}
\label{fig:edge-runtime-design}
\end{figure}

Các stage giao tiếp qua hàng đợi có kích thước giới hạn để tránh tăng latency vô
hạn. Khi consumer chậm, chính sách có thể bỏ frame cũ thay vì xử lý mọi frame.
Bộ đệm GPU được cấp phát khi load model và tái sử dụng trong suốt vòng đời
runtime. Metrics timer bao quanh từng stage và toàn pipeline.

Runtime nhận một cấu hình bất biến cho mỗi phiên chạy, gồm camera source, model
artifact, input shape, threshold, NMS và output sink. Thay đổi cấu hình tạo phiên
bản mới và được Edge Agent áp dụng có kiểm soát.